\begin{abstract}
\noindent\textbf{Introducción.} Las pérdidas no técnicas de energía se combaten con
inspecciones de campo de capacidad limitada, de modo que el problema práctico es decidir a
quién inspeccionar primero. La literatura reporta desempeños muy altos sin controlar que la
etiqueta registra fraude detectado y no fraude cometido.
\textbf{Objetivo.} Desarrollar un modelo de clasificación que ordene a los clientes según su
probabilidad de fraude a partir de su historial de facturación y sus atributos, para
priorizar inspecciones.
\textbf{Métodos.} Se usaron los datos públicos de la empresa tunecina STEG: 135.493 clientes
y 4.476.749 facturas, con 5,58\,\% de fraude detectado. El 20\,\% de los clientes se reservó
como conjunto de prueba evaluado una sola vez, y sobre el resto se aplicó validación cruzada
de cinco pliegues estratificada y agrupada por cliente. El historial de cada cliente se
censuró en su primer reemplazo de contador o cambio de código o de tarifa, por ser huella de
la inspección. Se compararon regresión logística, LightGBM y XGBoost con PR-AUC, ROC-AUC y
precisión en listas de tamaño fijo.
\textbf{Resultados.} Con el historial completo, las variables de su segunda mitad alcanzaron
un PR-AUC de 0,396 frente a 0,182 de la primera; con el historial censurado la diferencia
bajó a 0,032. El modelo seleccionado, LightGBM, obtuvo en el conjunto de prueba un PR-AUC de
0,191 (IC 95\,\% [0,174; 0,211]) y un ROC-AUC de 0,762, frente a 0,132 y 0,745 de una línea
base de exposición. En una lista del 1\,\% de los clientes encontró 118 fraudes, frente a 57
de esa línea base y 15 esperados al azar.
\textbf{Conclusiones.} El historial de facturación permite priorizar inspecciones con una
precisión 7,8 veces superior al azar. Evaluar sin retirar las huellas de la inspección
sobrestima el desempeño que un modelo tendría con clientes aún no inspeccionados.

\medskip\noindent\textbf{Palabras clave:} pérdidas no técnicas; detección de fraude;
priorización de inspecciones; fuga de información; \textit{gradient boosting}.
\end{abstract}
